\documentclass[10pt]{article}

% graphicx package, useful for including eps, pdf, and raster graphics
\usepackage{graphicx}
\DeclareGraphicsExtensions{.png,.png,.jpg,.pdf}

% basic packages
\usepackage{fancyhdr}
\usepackage{color}
\usepackage{parskip}
\usepackage{float}
\usepackage{microtype}
\usepackage{url}
\usepackage{hyperref}
\usepackage{booktabs}
\usepackage{makecell}
\usepackage{multirow}
\usepackage{pbox}
\usepackage{enumitem}
\usepackage{xeCJK}

% fonts
\usepackage{fontspec}
\setmainfont{DroidSans}[
  Extension   = .ttf,
  UprightFont = *,
  BoldFont    = *-Bold ]
\setmonofont{DroidSansMono}[
  Extension   = .ttf,
  UprightFont = * ]
\setCJKmainfont{NotoSansCJKsc}[
  Path        = /src/fonts/,
  Extension   = .otf,
  UprightFont = *-Regular,
  BoldFont    = *-Bold,
  Language    = Chinese Simplified ]

\hypersetup{colorlinks,urlcolor=blue}
\urlstyle{same}

% page layout
\usepackage[top=1in, bottom=1in, left=0.75in, right=0.75in, includefoot, heightrounded]{geometry}

% Useful for seeing the layout as you tweak it.
%\usepackage{showframe}

\setlength{\headheight}{1.1in}
\pagestyle{fancy}
\fancyhf{}

\lhead{\includegraphics[height=1in,clip]{logo}}
\rhead{\pbox[b][1in][c]{\textwidth}{
    \small
    \textbf{测试由以下实验室执行：}\\
    Northwest Genomics Center\\
    University of Washington School of Medicine\\
    Genome Sciences Box 355065\\
    3720 15th Ave NE\\
    Seattle, WA 98195-5065}}

\newcommand{\PageLine}{\rule{\textwidth}{0.25mm}}

\newcommand{\link}[2]{\href{#1}{#2}\footnote{#1}}

\begin{document}

\begin{center}
\Large
\textbf{NWGC SARS-CoV-2 (COVID-19) PCR 定性化验报告}
\end{center}

\bigskip

\begin{description}[font=\normalfont,align=left,labelwidth=8em]
\item [参与者姓名] \textbf{\VAR{pat_name|e}}
\item [出生日期] \textbf{\VAR{birth_date|e}}
\item [样本标识符] \textbf{\VAR{qrcode|e}}
\item [样本提交日期] \textbf{\VAR{collect_ts|e}}
\item [结果提供日期]
  %- if status_code == "never-tested"
  \textbf{没有进行检测}
  %- else
  \textbf{\VAR{result_ts|e}}
  %- endif
\item [样本类型] 鼻腔中部，自我采样
\end{description}

\PageLine

SARS-CoV-2 (COVID-19) PCR（Polymerase Chain Reaction，聚合酶链反应）定性结果

%- if status_code == "negative"
\textbf{未检测到}\\
%- elif status_code == "positive"
\textbf{检测到 SARS-CoV-2（阳性）}\\
%- elif status_code == "inconclusive"
\textbf{无定论}\\
%- elif status_code == "never-tested"
\textbf{无法检测}\\
%- else
\VAR{0/0} %# poor man's raise
%- endif

\PageLine

请注意：

\begin{itemize}
\item

  \textbf{此项目中所使用的样本和检测的类型，是金县及其他机构合作开展的公共卫生监测计划中的一部分。}
  您所接受的检测，并不能替代由医疗保健提供者提供的个人医疗护理。

\item

  \textbf{样本监管链：} 所检测的样本采用自我采集、自我贴标签的形式，并通过邮寄或配送服务提交给
  SCAN。其并非由持照医疗保健专业人员采集、贴标签和运输。这可能会影响检测结果，
  且无法排除假阴性的结果。后续的临床检测将更具保障。

\item

  \textbf{样本采集装置：} 中鼻甲骨拭子并不是FDA/CDC Emergency Use Authorization (EUA)
  （食品药品监督管理局/疾病控制和预防中心紧急使用授权）批准的样本类型。因此，
  未检测到/阴性结果可能无法排除 SARS-CoV-2 感染。

\item

  \textbf{Washington State Department of Health（华盛顿州卫生部）可能会对``阳性''结果进行验证。}
  如果您有任何疑问，应当与您的医疗保健提供者讨论此结果。

\item{
  \textbf{如果您的检测结果显示``未检测到''，这意味着未在我们的化验中检测到 SARS-CoV-2。}

  任何检测都不完美，您仍有较小的概率实际携带有 SARS-CoV-2 病毒但 COVID-19 检测结果为阴性。
  此检测并不采用由医疗保健提供者采集样本的标准方法，且样本是通过邮寄或配送服务提交的，
  而不是在医疗保健设施中采集和运送的。 无论您的检测结果如何，当您感觉身体不适时，
  最好的做法是待在家中并避免与其他人接触。如果您有其他健康问题，或是病情较重，
  或是对您的疾病感到担忧，您应当联系您的医疗保健提供者。 如果您的症状持续，
  您应当与您的医疗保健提供者讨论接受另一项 COVID-19 检测。
}

\item

  \textbf{如果您的检测结果显示``无定论''，} 这意味着我们实验室的 SARS-CoV-2
  靶向检测结果不确定。

  任何检测都不完美，无定论的结果仍有概率意味着您实际携带有导致 COVID-19 的 SARS-CoV-2 病毒。
  此检测并不采用由医疗保健提供者采集样本的标准方法，且样本是通过邮寄或配送服务提交的，
  而不是在医疗保健设施中采集和运送的。 无论您的检测结果如何，当您感觉身体不适时，
  最好的做法是待在家中并避免与其他人接触。如果您有其他健康问题，或是病情较重，
  或是对您的疾病感到担忧，您应当联系您的医疗保健提供者。 如果您的症状持续，
  您应当与您的医疗保健提供者讨论接受另一项 COVID-19 检测。

\item

  \textbf{如果您的检测结果显示``无法检测''，} 这意味着样本不足或是检测流程中出现了错误。

  由于检测无法进行，您仍有概率实际携带有导致 COVID-19 的 SARS-CoV-2 病毒。
  此检测并不采用由医疗保健提供者采集样本的标准方法，且样本是通过邮寄或配送服务提交的，
  而不是在医疗保健设施中采集和运送的。 无论您的检测结果如何，当您感觉身体不适时，
  最好的做法是待在家中并避免与其他人接触。如果您有其他健康问题，或是病情较重，
  或是对您的疾病感到担忧，您应当联系您的医疗保健提供者。 如果您的症状持续，
  您应当与您的医疗保健提供者讨论接受另一项 COVID-19 检测。

\item

  \textbf{阳性检测结果（显示为``检测到''或``检测到病毒''）将由检测结构报告给相关州和县卫生部。}
  您无需自行上报这些结果。 如果您的检测结果显示新型冠状病毒呈阳性，公共卫生部门将跟进您的情况，
  向您询问您的相关病况以及您可能接触的人。

\item

  \textbf{您可以将这些检测结果打印出来，或是通过电子邮件发给您的医疗保健提供者。}

\end{itemize}

\bigskip
\textbf{SARS-CoV-2 (COVID-19) PCR 定性化验方法：}

此化验旨在使用逆转录 PCR 来检测 SARS-CoV-2。此化验依靠的是扩增病毒基因组片段 Orf1B 和 S，
并采用内部人体 RNaseP 扩增对照。 University of Washington’s Northwest Genomics
Center（华盛顿大学西北基因组学中心）开发了这项检测并确定了其运行特性。
这项由实验室开发的检测已通过验证并已提交至 US Food and Drug Administration
（美国食品及药物管理局）以获取紧急使用授权。FDA 对于检测授权的独立审核为待处理状态。
此实验室经Clinical Laboratory Improvement Amendments（CLIA，临床实验室改进修正案）
认证为符合资格可进行高复杂度的临床实验室检测。

如需获取更多服务提供者信息，请参见 Washington State Department of Health 提供的以下链接：

\begin{itemize}
\item

  \link{https://www.doh.wa.gov/Portals/1/Documents/1600/coronavirus/COVID-SickExposed-ZH-Trd.pdf}{COVID-19
  确诊或疑似患者}

\item

  \link{https://www.doh.wa.gov/Portals/1/Documents/1600/coronavirus/COVIDexposed.pdf}{曾接触
  COVID-19 确诊病例的患者}

\item

  \link{https://www.doh.wa.gov/Portals/1/Documents/1600/coronavirus/COVIDconcerned.pdf}{有
  COVID-19 症状的无接触患者}

\end{itemize}

更多资源可在 Washington State Department of Health 和 CDC 网站上获取：

\begin{itemize}
\item

  \url{https://www.doh.wa.gov/Emergencies/NovelCoronavirusOutbreak2020COVID19/ChineseSimplified}

\item

  \url{https://chinese.cdc.gov/coronavirus/2019-ncov/index.html}

\end{itemize}

\end{document}
